% A height-bound-free method for extending integral point sets, with two applications
% J. Wilder, 2026-07-25
%
% Build:  pdflatex integral-heptagon-maximality.tex   (twice, for refs)
% Source of record: oracle/library/papers/INTEGRAL-HEPTAGON-MAXIMALITY-2026-07-25.md
% Every numerical value here is reproduced from the receipts in Section 7; none was typed by hand.
\documentclass[11pt,reqno]{amsart}

\usepackage[T1]{fontenc}
\usepackage[utf8]{inputenc}
\usepackage{amsmath,amssymb,amsthm}
\usepackage{array}
\usepackage{booktabs}
\usepackage[margin=1.15in]{geometry}
\usepackage[hidelinks]{hyperref}

\theoremstyle{plain}
\newtheorem{theorem}{Theorem}
\newtheorem{lemma}[theorem]{Lemma}
\newtheorem{proposition}[theorem]{Proposition}
\newtheorem{corollary}[theorem]{Corollary}

\numberwithin{equation}{section}

\newcommand{\dgp}{\dot{d}}
\newcommand{\Egp}{E_{\mathrm{gp}}}

\begin{document}

\title{A height-bound-free method for extending integral point sets,\\ with two applications}
\author{J. Wilder}
\date{\today}

\begin{abstract}
A planar \emph{integral point set} is a finite set of points with all pairwise Euclidean distances
integral; it is in \emph{general position} if no three of its points are collinear and no four
concyclic. Kreisel and Kurz~\cite{KK08} exhibited the first two seven-point examples, of diameters
$22270$ and $66810$, answering a question of Erd\H{o}s; whether an eight-point example exists is open.

We give an elementary method that decides, for any given integral point set $S$, whether a further
point lies at integral distance from every point of $S$ --- exhaustively and with no bound on how far
that point may lie. The observation is that for any point $X$ and any $P_i,P_j\in S$ the quantity
$c=\lvert XP_i\rvert-\lvert XP_j\rvert$ is an integer with $\lvert c\rvert\le d_{ij}$, a range that
does not depend on $\lvert X\rvert$; three such constraints determine $X$ by trilateration. The
resulting search is finite by construction, so no effective height bound is required.

Applying it we prove that both Kreisel--Kurz heptagons are \emph{maximal}: apart from their own
vertices, no point of the plane lies at integral distance from all seven vertices of either
configuration. Combining this with the exhaustive range of~\cite{KK08} gives the lower bound
$\dgp(2,8)>30000$ for the minimum diameter of an integral general-position octagon. We also record a
cardinality certificate --- if a triangle admits fewer than five admissible extension points then no
octagon contains it --- and verify it for $394$ triangles.
\end{abstract}

\maketitle

\section{Definitions and background}

Let $S\subset\mathbb{R}^2$ be finite. $S$ is an \emph{integral point set} if $\lvert PQ\rvert\in
\mathbb{Z}$ for all $P,Q\in S$, and is in \emph{general position} if no three points of $S$ are
collinear and no four are concyclic. Write $\dgp(2,n)$ for the minimum diameter of an $n$-point
integral point set in general position.

The known values are $\dgp(2,n)=1,8,73,174$ for $n=3,4,5,6$, and $\dgp(2,7)=22270$~\cite{KK08}.
Kreisel and Kurz further report that their diameter-$22270$ configuration is the \emph{only} such
heptagon of diameter at most $30000$, an exhaustive statement we use in Section~\ref{sec:maximality}.

By the Erd\H{o}s--Anning theorem~\cite{AE45} an integral point set that is not collinear is finite;
the usual proof is the hyperbola argument we make effective in Section~\ref{sec:method}. Every
non-collinear integral point set has a \emph{characteristic} $c$, the squarefree part of
$16\cdot\mathrm{Area}^2$ of any of its triangles, which is the same for all of them~\cite{Kem88};
both heptagons below have characteristic $2002=2\cdot 7\cdot 11\cdot 13$, and their points
accordingly lie in $\mathbb{Q}(\sqrt{2002})$.

Throughout, a point of $\mathbb{Q}(\sqrt{2002})^2$ is written $(x,y)$ meaning the real point
$(x,y\sqrt{2002})$.

\section{The method}\label{sec:method}

Let $S=\{P_0,\dots,P_{m-1}\}$ be an integral point set and let $X$ be any point with $\lvert
XP_i\rvert\in\mathbb{Z}$ for all $i$.

\begin{lemma}\label{lem:range}
For any $i\neq j$, the integer $c_j:=\lvert XP_0\rvert-\lvert XP_j\rvert$ satisfies $\lvert
c_j\rvert\le d_{0j}$.
\end{lemma}

\begin{proof}
The triangle inequality applied to $X,P_0,P_j$.
\end{proof}

The point of Lemma~\ref{lem:range} is that the range of $c_j$ is determined by $S$ alone and is
\emph{independent of how far $X$ lies from $S$} --- no height bound is needed, and none is assumed
anywhere below.

Fix a base index $0$ and two further indices $j,k$ with $P_0,P_j,P_k$ not collinear. Put
$r=\lvert XP_0\rvert$. Translating $P_0$ to the origin, write $u_i=x_i-x_0$, $w_i=y_i-y_0$,
$e_i=d_{0i}^2-c_i^2$ and $\Delta=u_jw_k-w_ju_k\neq 0$. Expanding $\lvert X-P_i\rvert^2=(r-c_i)^2$
against $\lvert X\rvert^2=r^2$ gives two equations linear in the coordinates of $X$:
\begin{align}
2u_j x + 2c\,w_j y &= e_j + 2rc_j,\\
2u_k x + 2c\,w_k y &= e_k + 2rc_k,
\end{align}
where $c$ is the characteristic. Their solution is affine in $r$:
\begin{equation}
x = A + Br,\qquad y = C + Dr,
\end{equation}
\begin{equation}
\begin{aligned}
A &= \frac{w_ke_j-w_je_k}{2\Delta}, &\qquad B &= \frac{w_kc_j-w_jc_k}{\Delta},\\[2pt]
C &= \frac{u_je_k-u_ke_j}{2c\Delta}, &\qquad D &= \frac{u_jc_k-u_kc_j}{c\Delta}.
\end{aligned}
\end{equation}
Substituting into $x^2+cy^2=r^2$ yields a quadratic in $r$ with rational coefficients:
\begin{equation}\label{eq:quadratic}
(B^2+cD^2-1)\,r^2 + 2(AB+cCD)\,r + (A^2+cC^2) = 0.
\end{equation}

\begin{theorem}\label{thm:method}
The set of points at integral distance from all of $P_0,P_j,P_k$ is finite and is enumerated exactly
by letting $(c_j,c_k)$ range over $[-d_{0j},d_{0j}]\times[-d_{0k},d_{0k}]$ and solving
\eqref{eq:quadratic} for $r$; each cell contributes at most two candidates.
\end{theorem}

This is Erd\H{o}s--Anning made effective: the grid is finite by Lemma~\ref{lem:range},
trilateration from three non-collinear points is injective, and every step is exact rational
arithmetic. Choosing the base triple to minimise $(2d_{0j}+1)(2d_{0k}+1)$ minimises the cost; for a
heptagon one takes the vertex incident to the two shortest edges.

\section{The two heptagons}\label{sec:heptagons}

Both configurations have characteristic $2002$. We restate them in the coordinates used here.

\medskip
\noindent\textbf{H1} (diameter $22270$):
\[
\begin{array}{l|rr}
 & x & y\\\hline
P_0 & 0 & 0\\
P_1 & 22270 & 0\\
P_2 & 26127018/2227 & 932064/2227\\
P_3 & 245363/17 & 3144/17\\
P_4 & 17615968/2227 & 238464/2227\\
P_5 & 56068/17 & 3144/17\\
P_6 & 19079044/2227 & -54168/2227
\end{array}
\]
\[
D(\mathrm{H1})=
\begin{pmatrix}
    0 & 22270 & 22098 & 16637 &  9248 &  8908 &  8636\\
22270 &     0 & 21488 & 11397 & 15138 & 20698 & 13746\\
22098 & 21488 &     0 & 10795 & 14450 & 13430 & 20066\\
16637 & 11397 & 10795 &     0 &  7395 & 11135 & 11049\\
 9248 & 15138 & 14450 &  7395 &     0 &  5780 &  5916\\
 8908 & 20698 & 13430 & 11135 &  5780 &     0 & 10744\\
 8636 & 13746 & 20066 & 11049 &  5916 & 10744 &     0
\end{pmatrix}
\]

\medskip
\noindent\textbf{H2} (diameter $66810$):
\[
\begin{array}{l|rr}
 & x & y\\\hline
P_0 & 0 & 0\\
P_1 & 66810 & 0\\
P_2 & 7690545/131 & 91800/131\\
P_3 & 78381054/2227 & 2796192/2227\\
P_4 & 98596712/2227 & 1148736/2227\\
P_5 & 91548738/2227 & 162504/2227\\
P_6 & 3314490/131 & 91800/131
\end{array}
\]
\[
D(\mathrm{H2})=
\begin{pmatrix}
    0 & 66810 & 66555 & 66294 & 49928 & 41238 & 40290\\
66810 &     0 & 32385 & 64464 & 32258 & 25908 & 52020\\
66555 & 32385 &     0 & 34191 & 16637 & 33147 & 33405\\
66294 & 64464 & 34191 &     0 & 34322 & 53244 & 26724\\
49928 & 32258 & 16637 & 34322 &     0 & 20066 & 20698\\
41238 & 25908 & 33147 & 53244 & 20066 &     0 & 32232\\
40290 & 52020 & 33405 & 26724 & 20698 & 32232 &     0
\end{pmatrix}
\]

Each was re-verified from its distance matrix alone: all $21$ distances integral and matching, no
three vertices collinear, no four concyclic.

\section{Maximality}\label{sec:maximality}

\begin{theorem}\label{thm:maximal}
Let $H$ be either heptagon of Section~\ref{sec:heptagons}. No point of the plane other than the seven
vertices of $H$ lies at integral distance from all seven vertices. In particular neither heptagon is
contained in an eight-point integral point set.
\end{theorem}

\begin{proof}
Apply Theorem~\ref{thm:method} with the base triple minimising the grid. Every cell was evaluated in
exact integer arithmetic; each candidate $r$ was accepted only when the quadratic vanished
identically over $\mathbb{Q}$, and the induced point was then tested for integrality of all seven
distances. A mandatory correctness gate required each run to recover all six non-base vertices, which
are themselves solutions; a run failing that gate was discarded rather than reported.

\begin{center}
\begin{tabular}{llrrrrcr}
\toprule
config. & base triple & grid cells & sol. & vertices & \textbf{non-trivial} & gate & time\\
\midrule
H1 & $(P_4;P_5,P_6)$ &   136{,}801{,}313 & 6 & 6 & \textbf{0} & 6/6 &  153\,s\\
H1 & $(P_5;P_4,P_0)$ &   205{,}982{,}337 & 6 & 6 & \textbf{0} & 6/6 &  223\,s\\
H2 & $(P_4;P_2,P_5)$ & 1{,}335{,}425{,}575 & 6 & 6 & \textbf{0} & 6/6 & 1530\,s\\
\bottomrule
\end{tabular}
\end{center}

Total: $1{,}678{,}209{,}225$ cells, no floating-point arithmetic anywhere in the decision path. H1 was
run twice from independent base triples, which agree.
\end{proof}

\begin{corollary}\label{cor:bound}
$\dgp(2,8)>30000$.
\end{corollary}

\begin{proof}
Suppose $O$ is an integral general-position $8$-point set of diameter at most $30000$. General
position is hereditary, so each of the eight $7$-point subsets of $O$ is an integral heptagon in
general position of diameter at most $30000$. By the exhaustive result of~\cite{KK08} every such
heptagon is H1. Then $O$ extends H1, contradicting Theorem~\ref{thm:maximal}.
\end{proof}

Corollary~\ref{cor:bound} depends on the exhaustive range reported in~\cite{KK08};
Theorem~\ref{thm:maximal} does not.

\section{A cardinality certificate}

For a triangle $T$ let $\Egp(T)$ denote the set of points at integral distance from all three vertices
of $T$ and in general position with $T$ (the general-position condition is necessary for each point of
an octagon individually, hence sound to impose pointwise). $\Egp(T)$ is finite and computable by
Theorem~\ref{thm:method}.

\begin{proposition}
If $\lvert\Egp(T)\rvert<5$ then no integral general-position $8$-point set contains $T$.
\end{proposition}

\begin{proof}
Such a set would consist of $T$ together with five further points, each lying in $\Egp(T)$.
\end{proof}

This is a proof rather than a failed search, and it is cheap: the count alone suffices. We computed
and independently verified it for $394$ triangles of characteristic $2002$. Verification re-enumerates
$\Egp(T)$ from the problem statement, sharing no code with the program that produced the certificates;
the independent counts agreed in every case, and the verifier correctly \emph{refuses} triangles with
$\lvert\Egp(T)\rvert\ge 5$, where the argument does not apply.

\section{Remarks}

The method is not specific to $n=7$: Theorem~\ref{thm:method} applies to any integral point set, and
the cost is set by the two shortest edges at the chosen base vertex, not by the diameter. Two
observations from the computations may be of independent interest.

First, the minimum edge of the known extremal configurations grows sharply with $n$: $1,5,26,68$ for
$n=3,4,5,6$, then $5780$ for H1 and $16637$ for H2. Any triangle of an integral heptagon therefore has
all edges at least $5780$, which excludes small seeds from consideration.

Second, among $63$ squarefree characteristics tested over $57{,}334$ seed triangles with sides at most
$260$, only $c=2002$ produced a six-point configuration; every other characteristic stopped at five or
fewer. The characteristic of both known heptagons is thus distinguished in this range, on $14\times$
fewer samples than the characteristics it beat.

\section{Reproduction}

\begin{center}
\begin{tabular}{ll}
\toprule
artifact & path\\
\midrule
method + closure driver & \texttt{oracle/kbk/engine/kk\_closure\_full\_exact.py}\\
heptagon reconstruction & \texttt{oracle/kbk/engine/kk\_heptagon.py}\\
closure receipts & \texttt{closure\_full\_exact\_\{H1,H1\_base540,H2\}.json}\\
cardinality certificates & \texttt{oracle/evidence/intdist-certificates/} (394 files)\\
independent verifier & \texttt{oracle.kernel.checkers.check\_intdist\_no\_octagon}\\
\bottomrule
\end{tabular}
\end{center}

\noindent Evidence receipt: \texttt{sha256
3c41ed8c959678eac586aa8ca316577a45bb284f0a2861b85bcfad5d7df189bd}.

\begin{thebibliography}{KK08}

\bibitem[AE45]{AE45}
N.~H. Anning and P.~Erd\H{o}s,
\emph{Integral distances},
Bull. Amer. Math. Soc. \textbf{51} (1945), 598--600.

\bibitem[Kem88]{Kem88}
A.~Kemnitz,
\emph{Punktmengen mit ganzzahligen Abst\"anden},
Habilitationsschrift, TU Braunschweig, 1988.

\bibitem[KK08]{KK08}
T.~Kreisel and S.~Kurz,
\emph{There are integral heptagons, no three points on a line, no four on a circle},
Discrete Comput. Geom. \textbf{39} (2008), 786--790.
\texttt{arXiv:0804.1303}.

\end{thebibliography}

\end{document}
